\section*{Bab \ref{ch:b2:himalaya-andes} --- Sabuk Tumbukan dan Sabuk Andes Dibandingkan}

\begin{solution}{exo:b2:himalaya-andes:1}
(a) $r = h\rho_c/(\rho_m - \rho_c) = 3.0\times2700/600 = \qty{13.5}{km}$. (b) $T = 35 + 3.0 +
13.5 = \qty{51.5}{km}$, atau $35 + 5.5\times3.0$. (c) $h = \Delta T(\rho_m - \rho_c)/\rho_m =
10\times600/3300 = \qty{1.8}{km}$.
\end{solution}

\begin{solution}{exo:b2:himalaya-andes:2}
(a) Palung depan Indo-Gangga; Sesar Dorong Frontal Utama; sub-Himalaya (Siwalik); Sesar Dorong Batas
Utama; Himalaya Kecil; Sesar Dorong Sentral Utama; Sekuen Himalaya Besar; Sistem Bidang Lepasan
Tibet Selatan; Sekuen Himalaya Tethys; lalu sutura Indus--Tsangpo dan kerak
Asia. (b) Sistem Bidang Lepasan Tibet Selatan. (c) Derajatnya naik ke utara: \omterm{def:b2:alpine-orogen:foreland}{molasse}
tak termetamorfosis, Himalaya Kecil berderajat menengah, Sekuen Himalaya Besar berderajat tinggi dan sebagian
meleleh; di atas STDS derajatnya turun tiba-tiba menjadi sedimen Tethys yang tak termetamorfosis atau
termetamorfosis lemah, karena sesar normal itu telah menyingkirkan batuan di antaranya.
\end{solution}

\begin{solution}{exo:b2:himalaya-andes:3}
(a) $196.5\ \text{m}/9200\ \text{yr} = \qty{0.02136}{m}$ per tahun, sekitar \qty{21.4}{mm} per tahun.
(b) Nilai itu \qty{20}{\percent} lebih besar daripada \qty{17.8}{mm} per tahun dari GPS, tetapi seukuran:
satu laju dirata-ratakan selama sembilan ribu tahun, yang lain selama dua puluh tahun.
(c) Hampir seluruh pemendekan di jajaran pegunungan itu diserap di front: MBT, MCT, dan
sesar-sesar lain di utara MFT, pada bujur ini, praktis tidak aktif; semuanya berakar
pada MHT yang sama, yang muncul di MFT.
\end{solution}

\begin{solution}{exo:b2:himalaya-andes:4}
(a) Tumbukan: sutura menandai samudra yang lenyap. (b) Tipe Andes: busur memerlukan slab
yang menunjam. (c) Tumbukan, secara khas: lelehan kerak dari kerak yang menebal dan terpanaskan (lelehan itu juga
terjadi dalam kerak Andes yang menebal, tetapi \omterm{def:b2:magma-emplacement:magma}{magma} busur mendominasi di sana). (d) Tipe Andes.
(e) Keduanya: Siwalik dan daratan depan Subandes. (f) Keduanya: MHT di Himalaya, bidang
antarmuka subduksi di Andes.
\end{solution}

\begin{solution}{exo:b2:himalaya-andes:5}
(a) $17.8/45 = \qty{40}{\percent}$ di timur, $20.5/45 = \qty{46}{\percent}$ di barat:
sekitar separuh. (b) Ke penebalan dan sesar-sesar Tibet dan bagian Asia lainnya, serta ke
ekstrusi blok-blok ke timur di sepanjang sesar-sesar mendatar besar. (c) $45 - 17.8 =
\qty{27.2}{mm}$ per tahun, atau \qty{272}{km} dalam \qty{10}{Ma}; $45 - 20.5 = \qty{24.5}{mm}$ per
tahun, atau \qty{245}{km}.
\end{solution}

\begin{solution}{exo:b2:himalaya-andes:6}
(a) $r = 4.5\times3.7 = \qty{16.65}{km}$; $\Delta F = 2700\times9.8\times3700\times(35\,000 +
1850 + 8325) = \qty{4.42e12}{N.m^{-1}}$. (b) Tibet: \qty{6.12e12}{N.m^{-1}}; Altiplano
memiliki \qty{72}{\percent} darinya. (c) Dengan $r = 4.5h$, $\Delta F = \rho_c g\,(Hh + 2.75h^2)$: untuk
$h \ll H$ suku $Hh$ dominan dan $\Delta F$ sebanding dengan $h$; untuk dataran tinggi yang tinggi,
suku dalam $h^2$ tumbuh lebih cepat. Setiap kilometer ketinggian berikutnya menuntut dorongan yang makin besar di
tepi-tepinya, sampai gaya yang tersedia tidak lagi dapat mengangkat dataran tinggi itu; dataran tinggi itu lalu menyebar
alih-alih naik.
\end{solution}

\begin{solution}{exo:b2:himalaya-andes:7}
(a) $2700\times9.8\times8848.86/2 = \qty{1.17e8}{Pa} = \qty{117}{MPa}$. (b) Pada
\qty{100}{MPa}: $0.85\times100 = \qty{85}{MPa}$; pada \qty{200}{MPa}: $0.85\times200 = 50 +
0.6\times200 = \qty{170}{MPa}$, kedua cabang itu sesuai. (c) $h = 2\sigma_d/(\rho g) =
2\times1.70\times10^8/(2700\times9.8) = \qty{12.8}{km}$: sebuah orde besaran, seorde dengan
puncak-puncak tertinggi, tidak jauh di atasnya.
\end{solution}

\begin{solution}{exo:b2:himalaya-andes:8}
(a) $11\times111.19\times\cos\ang{30} = \qty{1059}{km}$. (b) $9\times10^{-3}/1.059\times10^{6}
= \num{8.5e-9}$ per tahun (antara \num{6.6e-9} dan \num{10.4e-9} untuk 7 dan \qty{11}{mm} per tahun),
yaitu $8.5\times10^{-9}/3.156\times10^{7} = \qty{2.7e-16}{s^{-1}}$. (c) Kedua rentang itu
bertumpang-tindih: $(12.4 \pm 4.7) \times 10^{-9}$ mencakup 7.7 sampai \num{17.1e-9}. Dataran tinggi bagian selatan di utara
Nepal meregang secepat, atau sedikit lebih cepat daripada, sabuk itu secara keseluruhan.
\end{solution}

\begin{solution}{exo:b2:himalaya-andes:9}
(a) $0.9/6.6 = \qty{14}{\percent}$. (b) $M_0 = 10^{1.5\times8.5 + 9.1} = 10^{21.85} =
\qty{7.1e21}{N.m}$. (c) $7.1\times10^{21}/6.6\times10^{19} = 107$: satu gempa bumi semacam itu sekitar
setiap 110 tahun di bawah Nepal, jauh lebih banyak daripada yang dicatat sejarah; entah gempa bumi yang jauh lebih besar
akan datang, atau sebagian defisit dilepaskan dengan cara lain.
\end{solution}

\begin{solution}{exo:b2:himalaya-andes:10}
(a) Ambillah kedalaman kompensasi di dasar kerak Altiplano, $61.5 - 3.7 =
\qty{57.8}{km}$ di bawah muka laut. Kerak Puna mencapai $50.5 - 4.4 = \qty{46.1}{km}$, sehingga
kolom Puna mengandung \qty{11.7}{km} mantel di atas kedalaman itu: $2800\times61.5 =
\rho_P\times50.5 + 3300\times11.7$. (b) $\rho_P = (172\,200 - 38\,610)/50.5 =
\qty{2645}{kg.m^{-3}}$. (c) Selisihnya, \qty{155}{kg.m^{-3}}, terletak di ujung bawah
rentang 150 sampai \qty{200}{kg.m^{-3}} yang dikutip: kerak saja akan memerlukan kerak Puna yang luar biasa ringan.
Kolom yang lebih ringan di bawahnya lebih mudah menjelaskannya: litosfer mantel yang tersingkir oleh delaminasi,
yang digantikan oleh astenosfer yang panas dan ringan.
\end{solution}

\begin{solution}{exo:b2:himalaya-andes:11}
(a) $g = 105/\tan\ang{30} = \qty{181.9}{km}$. (b) $g = 90/\tan\ang{30} + 15/\tan\ang{4} =
155.9 + 214.5 = \qty{370.4}{km}$: busur berpindah \qty{188.5}{km} ke pedalaman. (c) Slab itu tidak pernah
mencapai \qty{105}{km} di bawah lempeng sepanjang \qty{300}{km} itu, dan tidak ada baji mantel panas
di atasnya: tidak ada \omterm{def:b2:magma-emplacement:magma}{magma} yang terbentuk dan busur itu mati sampai slab tenggelam lagi, lebih jauh ke timur.
\end{solution}

\begin{solution}{exo:b2:himalaya-andes:12}
(a) \qty{21}{mm} per tahun sama dengan \qty{21}{km} per juta tahun: $21\times37.5 =
\qty{787.5}{km^3}$ per kilometer jajaran pegunungan per juta tahun. (b) $300/787.5 =
\qty{38}{\percent}$. (c) Dalam penebalan jajaran pegunungan dan Tibet selatan, dalam kerak yang
terbawa turun bersama India di bawah Tibet, dan dalam materi yang diteruskan MHT ke utara.
\end{solution}

\begin{solution}{pb:b2:himalaya-andes:1}
\textbf{1.} Tujuh sampel dari \ang{28.0} sampai \ang{37.0}~LU: $34\,550/7 = \qty{4936}{m}$, $h_T
= \qty{4.9}{km}$.
\textbf{2.} Lima sampel dari \ang{-69} sampai \ang{-65}: $18\,650/5 = \qty{3730}{m}$, $h_A =
\qty{3.7}{km}$, nilai yang dikutip untuk Altiplano; Puna, lebih ke selatan, lebih tinggi.
\textbf{3.} Tibet: sekurang-kurangnya $9\times111.19 = \qty{1001}{km}$; Altiplano: sekurang-kurangnya $4\times
104.5 = \qty{418}{km}$.
\textbf{4.} $(5000 - 70)/(1.5\times111\,190) = 0.0296$: \qty{3.0}{\percent}.
\textbf{5.} Profil Andes dimulai di dasar samudra dan di parit, beberapa kilometer di bawah
muka laut, di barat dataran tinggi; profil Tibet dimulai di dataran benua dekat muka laut,
dan berakhir di cekungan Tarim yang tinggi.
\textbf{6.} $r = 4.5\times4.9 = \qty{22.05}{km}$; $T_T = 35 + 4.9 + 22.05 = \qty{62.0}{km}$.
\textbf{7.} $r = 4.5\times3.7 = \qty{16.65}{km}$; $T_A = 35 + 3.7 + 16.65 = \qty{55.4}{km}$.
\textbf{8.} Tibet: \qty{62.0}{km} dibandingkan dengan 70 sampai \qty{75}{km}, \qty{10.5}{km} kurang dari nilai
tengah \qty{72.5}{km} (\qty{14}{\percent}); Altiplano: \qty{55.4}{km} dibandingkan dengan 59 sampai
\qty{64}{km}, \qty{6.1}{km} kurang dari \qty{61.5}{km} (\qty{10}{\percent}). Tibet meleset lebih jauh.
\textbf{9.} $72.5 - 35 = 4.9\times3300/(3300 - \rho_c)$: $3300 - \rho_c = 16\,170/37.5 =
431.2$, $\rho_c = \qty{2869}{kg.m^{-3}}$.
\textbf{10.} $26.5 = 3.7\times3300/(3300 - \rho_c)$: $3300 - \rho_c = 460.8$, $\rho_c =
\qty{2839}{kg.m^{-3}}$.
\textbf{11.} Sebagian beban dipikul oleh pelengkungan lempeng, bukan oleh akar lokal;
kerak normal mungkin tidak setebal \qty{35}{km}; mantel di bawahnya mungkin lebih ringan atau lebih rapat
daripada yang diandaikan; Moho seismik mungkin bukan dasar kerak dalam arti massa jenis.
\textbf{12.} Dasar kerak Tibet terletak pada $72.5 - 4.9 = \qty{67.6}{km}$, dasar kerak India pada
\qty{37.5}{km}: $t = \qty{30.1}{km}$. $\Delta g = 2\pi\times6.674\times10^{-11}\times370\times
30\,100 = \qty{4.67e-3}{m.s^{-2}}$, yaitu \qty{-467}{mGal} relatif terhadap India, dekat dengan
\qty{460}{mGal} yang teramati.
\textbf{13.} $W_0 = 440\times60/35 = \qty{754.3}{km}$; $S = \qty{314.3}{km}$.
\textbf{14.} $(314.3 - 250)/314.3 = \qty{20}{\percent}$, kekurangan yang didapati di Andes
Tengah.
\textbf{15.} $0.55\times314.3 = \qty{173}{km}$, lebih kecil daripada \qty{250}{km} dari \omterm{def:b2:alpine-orogen:balanced}{penampang
seimbang}, yang dengan demikian akan lebih dari cukup.
\textbf{16.} Penutup yang memendek di atas \omterm{def:b2:alpine-orogen:decollement}{décollement} harus diimbangi oleh pemendekan
batuan dasar di bawahnya, lebih jauh ke barat, karena keduanya harus kehilangan panjang yang sama di seluruh sabuk; pemendekan batuan dasar
itu menebalkan kerak di bawah Cordillera Timur dan dataran tinggi.
\textbf{17.} $314.3/30 = \qty{10.5}{mm}$ per tahun, $10.5/72 = \qty{15}{\percent}$
konvergensi.
\textbf{18.} Sisa itu habis oleh subduksi: lempeng Nazca tenggelam di bawah benua.
\textbf{19.} $r = \qty{22.05}{km}$; $\Delta F = 2700\times9.8\times4900\times(35\,000 + 2450
+ 11\,025) = \qty{6.28e12}{N.m^{-1}}$.
\textbf{20.} $2700\times9.8\times3700\times45\,175 = \qty{4.42e12}{N.m^{-1}}$; rasionya
0.70.
\textbf{21.} $6.28\times10^{12}/10^{5} = \qty{6.3e7}{Pa}$, sekitar \qty{63}{MPa}.
\textbf{22.} $9\times10^{-3}/1.06\times10^{6} = \num{8.5e-9}$ per tahun.
\textbf{23.} Dorongan India bekerja dalam arah utara--selatan dan diimbangi, di dataran tinggi itu, oleh
beban dataran tinggi itu sendiri: dalam arah utara--selatan dataran tinggi itu tidak dapat lagi menebal, sehingga tegangan vertikal,
yaitu bebannya, adalah yang terbesar; tegangan terkecil berarah timur--barat, ke arah tepi-tepi bebas Asia,
dan dataran tinggi itu mengalami ekstensi ke arah itu, pada celah-celah berarah utara--selatan.
\textbf{24.} \textbf{\qty[text-series-to-math=true]{62}{km}}, sekitar \qty{86}{\percent}
dari \qty{72.5}{km} yang terukur.
\end{solution}
